%% ============================================================
%%  Chapitre 6 — Patterns d'agents avancés
%% ============================================================
\chapter{Patterns d'agents avancés}
\label{ch:patterns}

Ce chapitre détaille les trois patterns de raisonnement avancés intégrés
dans \textit{Scientific Watch Agent} : le pattern ReAct pour l'expansion
de requêtes, le pattern Reflexion pour la révision de synthèse, et
ProTeGi pour l'optimisation automatique des prompts.

\section{Pattern ReAct : le QueryExpander}
\label{sec:react}

\subsection{Motivation}

Un utilisateur formule son besoin de recherche en termes naturels :
``\textit{fake news detection}''. Cette requête brute, soumise directement
à OpenAlex, peut manquer des travaux importants formulés différemment
(``\textit{misinformation}'' ou ``\textit{rumor detection}'').
L'expansion de requêtes consiste à générer automatiquement plusieurs
reformulations diversifiées du topic, de façon à maximiser la couverture
de la littérature pertinente.

Le pattern ReAct \cite{yao2022react} est particulièrement adapté à cette
tâche car il permet au LLM de \textit{tester} ses hypothèses en sondant
réellement la base de données avant de formuler la requête suivante,
évitant ainsi les expansions basées sur des présupposés erronés.

\subsection{Boucle ReAct dans le QueryExpander}

La boucle ReAct du \texttt{QueryExpander} fonctionne comme suit :

\begin{algorithm}[htbp]
  \caption{Boucle ReAct du QueryExpander}
  \label{alg:react}
  \begin{algorithmic}[1]
    \State $\text{queries} \gets [\text{topic}]$  \Comment{Requête initiale forcée}
    \State $i \gets 0$
    \While{$i < \text{MAX\_ITERATIONS}$ \textbf{and} $\text{stop\_reason} \neq \text{"done"}$}
      \State $\text{thought, query, stop} \gets \text{LLM}(\text{topic, history, observations})$
      \State $\text{count, concepts, examples} \gets \text{OpenAlex.probe}(\text{query})$
      \State $\text{observations} \mathrel{+}= (\text{query, count, concepts, examples})$
      \If{$\text{query} \notin \text{queries}$}
        \State $\text{queries} \mathrel{+}= [\text{query}]$
      \EndIf
      \State $i \mathrel{+}= 1$
    \EndWhile
    \State \Return $\text{queries}$
  \end{algorithmic}
\end{algorithm}

À chaque itération, l'agent LLM reçoit le contexte complet (topic initial,
historique des pensées, résultats des sondes) et produit une sortie
structurée \texttt{ReActThoughtAction} comprenant :
\begin{itemize}
  \item \texttt{thought} : son raisonnement sur la requête à formuler ;
  \item \texttt{search\_query} : la requête concrète pour OpenAlex ;
  \item \texttt{stop\_reason} : ``continue'' ou ``done'' (si l'exploration
    est jugée suffisante).
\end{itemize}

\paragraph{Contrainte sur la première itération.}
La requête de la première itération est forcée à être le topic original
exact, garantissant que le corpus inclura toujours les articles les plus
directement pertinents avant d'explorer des reformulations.

\subsection{Observation : le feedback de l'environnement}

L'observation renvoyée à l'agent après chaque sonde contient :
\begin{itemize}
  \item Le nombre d'articles retournés par OpenAlex ;
  \item Les cinq concepts les plus associés aux résultats ;
  \item Les titres des trois premiers articles.
\end{itemize}

Ces informations permettent à l'agent d'évaluer la qualité de sa requête
et d'ajuster sa stratégie. Par exemple, si une requête retourne moins de
10 articles, l'agent peut choisir d'élargir la formulation. Si les concepts
associés ne correspondent pas au domaine visé, il peut reformuler en
ajoutant des termes discriminants.

\section{Pattern Reflexion : le couple Synthesizer--Critic}
\label{sec:reflexion}

\subsection{Motivation}

La synthèse d'un corpus de publications est une tâche complexe, sujette
à des défauts difficiles à détecter automatiquement : omissions d'approches
importantes, généralisations abusives, formulations vagues, ou inconsistances
entre sections. Le pattern Reflexion \cite{shinn2023reflexion} adresse ce
problème en introduisant un agent critique qui évalue la synthèse et
formule des retours verbaux que le synthétiseur peut intégrer dans une
révision.

\subsection{Cycle Reflexion}

Le cycle Reflexion implique deux agents en boucle :

\begin{enumerate}
  \item Le \textbf{Synthesizer} génère ou révise la synthèse du domaine
    (\texttt{Synthesis}) en prenant en compte le retour du \texttt{Critic}
    du cycle précédent (si disponible) ;
  \item Le \textbf{Critic} évalue la synthèse reçue selon quatre axes
    et produit un retour structuré \texttt{CriticFeedback} ;
  \item Si \texttt{needs\_revision = True} et que le nombre d'itérations
    est inférieur à \texttt{MAX\_REFLEXION\_ITERATIONS = 3}, l'arête
    conditionnelle de LangGraph renvoie vers le \texttt{Synthesizer}.
\end{enumerate}

\subsection{Axes d'évaluation du Critic}

Le \texttt{Critic} évalue la synthèse sur quatre axes, chacun noté de
1 à 5 par le LLM :

\begin{description}
  \item[Fidelity (fidélité).] La synthèse est-elle fidèle au contenu
    des articles résumés ? Ne contient-elle pas d'affirmations non
    supportées par les sources ?
  \item[Completeness (exhaustivité).] La synthèse couvre-t-elle les
    principales approches et contributions du domaine, sans omission
    majeure ?
  \item[Specificity (spécificité).] La synthèse est-elle suffisamment
    précise, avec des noms de méthodes, des chiffres et des exemples
    concrets ? Évite-t-elle les généralités vagues ?
  \item[Consistency (cohérence).] Les différentes sections de la synthèse
    sont-elles cohérentes entre elles ? Y a-t-il des contradictions ?
\end{description}

Une révision est déclenchée si la qualité globale est ``poor'' (score
moyen $< 3$) ou ``acceptable'' (score moyen $< 4$).
La qualité ``good'' (score $\geq 4$) arrête le cycle.

\subsection{Mécanisme de sécurité}

En cas d'échec de l'appel LLM du \texttt{Critic} (timeout, erreur API),
le système approuve automatiquement la synthèse courante (\texttt{force\_approve}),
permettant au pipeline de continuer sans bloquer. Ce mécanisme garantit
la robustesse du système en production.

\subsection{Exemple d'itération Reflexion}

Le tableau~\ref{tab:reflexion_example} illustre l'évolution qualitative
d'une synthèse sur le domaine ``anomaly detection'' après une itération
Reflexion.

\begin{table}[htbp]
  \centering
  \caption{Exemple de retour Critic et amélioration après Reflexion.}
  \label{tab:reflexion_example}
  \begin{tabular}{p{0.45\textwidth}p{0.45\textwidth}}
    \toprule
    \textbf{Synthèse initiale (iter. 1)} & \textbf{Retour Critic} \\
    \midrule
    ``Les méthodes de détection d'anomalies utilisent des approches
    diverses incluant l'apprentissage profond et les méthodes statistiques.''
    &
    ``Specificity faible (2/5) : aucune méthode nommée.
    Completeness insuffisante (2/5) : autoencoders, Isolation Forest,
    VAE non mentionnés.'' \\
    \bottomrule
    \multicolumn{2}{l}{\textbf{Synthèse révisée (iter. 2)}} \\
    \multicolumn{2}{p{0.92\textwidth}}{``Les approches dominantes incluent
    les autoencoders et VAE (Variational Autoencoders) pour la détection
    non supervisée, Isolation Forest pour la détection rapide à haute dimension,
    et LSTN-AD pour les séries temporelles. Les benchmarks MVTec AD et
    KDD Cup 1999 sont les plus utilisés.''} \\
  \end{tabular}
\end{table}

\section{ProTeGi : optimisation automatique des prompts}
\label{sec:protegi}

\subsection{Motivation}

Le \texttt{Summarizer} utilise un prompt qui décrit la tâche de résumé
au LLM. La qualité de ce prompt affecte directement la qualité des résumés
produits, et donc la qualité de l'ensemble du pipeline. Une optimisation
manuelle est possible mais coûteuse et dépendante du modèle utilisé.
ProTeGi \cite{pryzant2023automatic} permet d'automatiser cette optimisation.

\subsection{Fonctionnement de ProTeGi}

ProTeGi adapte le concept de descente de gradient au domaine textuel.
L'algorithme procède par itérations :

\begin{enumerate}
  \item \textbf{Évaluation} : le prompt courant est utilisé pour résumer
    un batch d'articles ; les résumés sont évalués par rapport à des
    références humaines via ROUGE-L \cite{lin2004rouge} et BERTScore
    \cite{zhang2020bertscore} ;
  \item \textbf{Gradient textuel} : un LLM ``critique'' analyse les
    erreurs des résumés produits et formule des suggestions d'amélioration
    du prompt ;
  \item \textbf{Mise à jour} : les meilleures suggestions sont incorporées
    dans une nouvelle version du prompt ;
  \item \textbf{Sélection} : le meilleur prompt parmi les candidats est
    retenu pour l'itération suivante.
\end{enumerate}

\subsection{Résultats expérimentaux}

Nous avons conduit 14 runs d'optimisation ProTeGi sur 6 configurations
de LLM différentes (claude-haiku + claude-sonnet, gpt-4o-mini + gpt-4o,
llama-3.1-8b + gpt-4o, llama-3.1-8b + claude-sonnet, deepseek-v3.2,
nebius-qwen3), en modes structuré et narratif.

Le tableau~\ref{tab:protegi_results} présente les résultats obtenus
en mode narratif, qui s'est avéré le plus bénéficié par l'optimisation.

\begin{table}[htbp]
  \centering
  \caption{Résultats ProTeGi (mode narratif) — ROUGE-L et BERTScore.}
  \label{tab:protegi_results}
  \rowcolors{2}{lightgray}{white}
  \begin{tabular}{lccc}
    \toprule
    \textbf{Configuration} & \textbf{ROUGE-L init.} & \textbf{ROUGE-L final} & \textbf{$\Delta$} \\
    \midrule
    haiku + sonnet & 0.31 & 0.38 & $+22\,\%$ \\
    gpt-4o-mini + gpt-4o & 0.29 & 0.35 & $+21\,\%$ \\
    llama-8b + gpt-4o & 0.27 & 0.33 & $+22\,\%$ \\
    deepseek-v3.2 & 0.33 & 0.37 & $+12\,\%$ \\
    nebius-qwen3 & 0.30 & 0.36 & $+20\,\%$ \\
    \bottomrule
  \end{tabular}
\end{table}

L'amélioration moyenne en ROUGE-L est de $+19\,\%$ sur les 5 configurations
testées. Le gain est plus modeste en mode structuré ($+8\,\%$ en moyenne),
ce format étant plus contraint et donc moins sensible à la formulation
du prompt.

\paragraph{Analyse qualitative.}
L'examen des prompts optimisés révèle que ProTeGi a systématiquement
ajouté des instructions explicites sur : la longueur cible du résumé
(``\textit{entre 150 et 250 mots}''), la nécessité de mentionner les
performances numériques (``\textit{inclure les métriques quantitatives
reportées}''), et l'interdiction des phrases vagues (``\textit{éviter
les généralisations sans exemples concrets}'').

\section{Discussion}

Les trois patterns présentés dans ce chapitre jouent des rôles complémentaires
dans la qualité globale du système.

ReAct améliore la \textbf{couverture} en diversifiant les requêtes de recherche,
réduisant le risque de manquer des travaux importants formulés différemment.
Reflexion améliore la \textbf{précision} de la synthèse en permettant une
révision guidée par un retour structuré. ProTeGi améliore la \textbf{qualité
de présentation} des résumés en optimisant les instructions données au LLM.

Ces trois améliorations sont orthogonales et cumulatives : un système intégrant
les trois patterns bénéficie simultanément d'une meilleure couverture,
d'une meilleure précision et d'une meilleure qualité de présentation.
